% $Id: ToR.tex,v 1.64 2025/07/26 08:48:56 brandenb Exp $
\documentclass{article}
\setlength{\textwidth}{168mm}
\setlength{\oddsidemargin}{-4mm}
%\setlength{\textheight}{230mm}
\setlength{\textheight}{245mm} %(AB: to make it fit on 2 pages)
\setlength{\topmargin}{-22mm}
\usepackage{graphicx,natbib,bm,url,html,color}
%\usepackage{graphicx,natbib,bm,url,color}
\thispagestyle{empty}
\def\red{\textcolor{red}}
\def\blue{\textcolor{blue}}
\title{Pencil Code Steering Committee (PCSC)}
\author{Terms of Reference}
\date{\small\today,~ $ $Revision: 1.64 $ $}
\begin{document}
\maketitle

\section{Purpose}

The PCSC is a standing committee whose purpose is to
\begin{itemize}\itemsep-2pt
\item
review and, when necessary, revise the Terms of Reference
(ToR, this document),

\item
review and be responsible for the public information about the code
on the {\sc Pencil Code} websites,

\item
help initiating the continuation of annual {\sc Pencil Code} User Meetings,

\item
invite Owners (=administrators in pencil-code.org)
and decide about their privileges and responsibilities,

\item
assess and check that the list of Committers and Owners is up-to-date\footnote{
\url{../committers/list.pdf}},

\item
deal with general concerns to the {\sc Pencil Code} community
(major code changes, conflicts, complaints),

\item
decide about major discussion items that need to be deferred to the next
Pencil Code User Meeting,

\item
organize elections of the PCSC (every 3 years).

\end{itemize}

\section{Organization of the PCSC}

\subsection{Chair}

The PCSC elects its chair for the duration of 1--2 years.
The chair continues
until a new chair has been elected or re-elected, which will
be the case prior to the first meeting after the chair's end of term.
Rotation of the chair is encouraged.
The chair selects a secretary for each meeting.
Except for changes to the ToR, or if a veto is raised, the quorum of meetings is 3/5.

\subsection{Members}

The PCSC has five members.
They are elected by the group of voters that comprises,
all {\em Owners of the Pencil Code}\footnote{
Owners of the {\sc Pencil Code} (as of \today):
\newpage
\begin{tabular}{cllll}
& Name & UserID & Affiliation & Country \\
\hline
& Sven Bingert       & bingert         & GWDG (G\"ottingen)   & Germany \\
$\ast$ & Philippe Bourdin   & PABourdin       & IWF (Graz)           & Austria \\
$\ast$ & Axel Brandenburg&AxelBrandenburg&Nordita             &Sweden\\
& Simon Candelaresi  & iomsn           & University of Glasgow & UK \\% (*) (2019-12-07)
$\ast$ & Piyali Chatterjee & mppiyali  & Indian Inst Astrophys & India\\%(2019-07-31)
& Boris Dintrans     & bdintrans       & IRAP Toulouse        & France \\
& Wolfgang Dobler    & wdobler         & Potsdam/Berlin       & Germany \\ % (*)
& Frederick Gent     & gent            & Aalto Univ.          & Finland \\%(2019-12-07)
$\ast$ & Nils Haugen & nilshau         & SINTEF (Trondheim)   & Norway \\
& Tobias Heinemann   & tobson          & NBIA (Copenhagen)    & Denmark \\
& Anders Johansen    & ajohan          & Lund Univ (Lund)     & Sweden \\
& Maarit Korpi-Lagg  & mjkorpi         & Aalto Univ.          & Finland \\ % (*)
& Petri K\"apyl\"a   & pkapyla         & Freiburg             & Germany \\ % (*)
& Wladimir Lyra      & wlyra           & Univ New Mexico      & USA \\
& Dhrubaditya Mitra  & dhrubaditya     & Nordita (Stockholm)  & Sweden \\ % (*)
$\ast$ & Matthias Rheinhardt& mrheinhardt & Aalto Univ (Espoo)& Finland \\
& Alberto Roper Pol  &                 & Geneva University    & Switzerland \\
& Jennifer Schober   & JenSchober      & EPFL (Lausanne)      & Switzerland \\
& Joern Warnecke     & jwarne          & MPS (G\"ottingen)    & Germany \\ % (*)
& Chao-Chin Yang     & chaochinyang    & Univ Nevada (Las Vegas) & USA \\ % (*)
& Hongzhe Zhou       & zhouhz1992      & TDLI Institute, Shanghai & China \\ % (*) (2022-09-01)
%latex comments: joining date, non-avail for PCSC
\end{tabular}
\newpage
Asterisks denote PCSC members (since November 2020).
Its current chair (from June 28, 2023) is Philippe Bourdin.
}
and those Committers who have participated in at least two
{\sc Pencil Code} User Meetings, and have made at least 20 check-ins to
the code within the last 5 years of the election\footnote{Use {\tt
git log --since 'today minus 5 years' --pretty=format:'\%an' | sort| uniq -c | sort -rn} \\
or select the time interval on github's graphics interface}.
Eligible members of the PCSC are Owners of the {\sc Pencil Code}
who have not declared their unavailability prior to the elections.
The vote is anonymous and voting will be carried out electronically.
Voting will be prepared by one of the Owners and approved by an observer.
Both will be nominated by the chair of the PCSC.

If the outcome of the vote is ambiguous, a runoff round with the candidates who have
equal vote counts is held.
If the runoff round ends ambiguously again, the observer decides freely
upon further measures to avoid runoff rounds in long sequence.
These might include drawing lots or a willful decision of the observer.

\section{Responsibilities of the PCSC}
The PCSC reports to the Owners with a critical self-assessment.
It is responsible for deciding what are major issues to be discussed.
The PCSC is also responsible for monitoring the long-term integrity
of the {\sc Pencil Code} through the help of the auto-test.
It also decides about Owners privileges.
Owners, in turn, have the responsibility to maintain, improve,
and disseminate the code.
No individual Owner is allowed to delete it.

\subsection{Meeting \& Reporting}

The PCSC convenes via telecon when issues emerge and
at least once a year at least 1--2 months before the next
Pencil Code User Meeting.
The PCSC reports via email to the Owners.
Upcoming meetings are announced to the Owners
to solicit additional agenda items.
The minutes will be available online.

\subsection{Recurrent tasks}

The PCSC reviews major code change requests or suggestions.
(A major code change is one that cannot be reverted or avoided
by switching modules, for example.)
Owners may request that a particular item be discussed and thereby
request that a PCSC meeting be held.
A major code change may also be one that is sufficiently complex
to foresee resulting major side effects.
The PCSC is considering to give voting rights to individual Committers
of the {\sc Pencil Code}.

\subsection{Review}

The PCSC reviews its operations at each Pencil Code User Meeting.
It is responsible for making changes to the ToR.
All major changes must be approved by all Owners with a 2/3 majority.
The PCSC also oversees the list of new Committers added by the Owners and makes
sure the full names and affiliations of all Committers are known.

\subsection{Complaints}

Complaints\footnote{including violations of the
\url{https://github.com/pencil-code/pencil-code/blob/master/license/CODE_OF_CONDUCT.md}
(CoC). The CoC specifies the Community Leaders, i.e., the Owners, as
responsible for its enforcement.} relevant to the Pencil Code community should be
directed to all 5 members of the PCSC.
When needed, the PCSC may then assemble an ad hoc committee
to resolve the problem.
The ad hoc committee, which decides by simple majority,
should consist of five members from the Pencil Code community.
If personal accusations are involved, both the accuser and the
accused are allowed to appoint one member each who is not part of the conflict.

\vfill\bigskip\noindent\tiny\begin{verbatim}
$Header: /var/cvs/brandenb/public_html/pencil-code/PCSC/ToR/ToR.tex,v 1.64 2025/07/26 08:48:56 brandenb Exp $
\end{verbatim}

\end{document}
